\documentclass{article}
\usepackage{graphicx}
\usepackage{braket}
\usepackage{amsmath}
\usepackage{amssymb}
\usepackage[T1]{fontenc}
\usepackage{lmodern}
\usepackage[english]{babel}
\usepackage[a4paper,top=2.2cm,bottom=2.2cm,left=2.3cm,right=2.3cm]{geometry}
\usepackage{microtype}
\usepackage[hidelinks]{hyperref}

\title{CQG paper II -- TK response 1}
\author{Iman Rosignoli}
\date{August 31, 2026}

\begin{document}

\maketitle

My thanks to both referees.

The first report was notably specific: it found defects I had not managed to
detect, one of which turned out to be more serious than the report could establish
from outside, and correcting it led me to a stronger result than the one I had
obtained initially. The second report allowed me to go deeper into the geometric
language surrounding my problem and to envisage a possible experimental test of the
work; the paper is now the richer for it.

The revision follows the first referee's list of points, in the order given.

\begin{enumerate}
    \item The exterior/contact/analytic-continuation triple is now made manifest
    throughout the work. Section~2.1 proves, as Lemma~2, the compactness criterion
    the referee outlined, from which it follows that the compact control problem
    terminates at \(r=2M\) for \(W=\partial_\eta\); Protocol~1 fixes the
    labelling (E)/(C)/(A) and is applied throughout the paper, figures included.

    \item The extremal-versus-minimiser statements are corrected in both papers,
    and the claim that numerical perturbations establish minimality is removed.

    \item The \(J=-J_c\) contradiction is resolved by rederiving the result. Both
    printed statements were wrong; the correct result is a four-case
    classification, Proposition~9, and all the material involved has been
    regenerated.

    \item The archive is rebuilt: the missing module has been restored, slopes and
    checksums are now generated rather than transcribed, and a manifest of twenty
    entries associates the artefacts with commands and checksums.
\end{enumerate}

The exterior retrograde separatrix is now the principal exact result, as
recommended, and the title no longer speaks of a trichotomy.

\section{Referee 1}

\subsection{Major comment 1. The control domain terminates at the conformal
stationary limit.}
\textbf{Reply.} Fully accepted; this is the structural change of the revision.

Lemma~2 states that the rail slice is nonempty and compact if and only if \(W\) is
timelike and \(\hat E \geq |W|\), degenerating to a single point at
\(\hat E = |W|\), and is empty or noncompact otherwise. From the form of
\(g(\partial_\eta,\partial_\eta)=-A^2(1-2M/r)\) one deduces that the two conditions
are \(r>2M\) and \(\bar v^2 = 1 - A^2f/\hat E^2 >0\).

The text states explicitly that a change of coordinates does not alter this
conclusion. Doran coordinates regularise the formulae at the points where
\(\Delta = 0\), but leave the inner product \(g(W,W)\) unchanged. A selector
anchored to Doran, or to ZAMO, is a different problem and is not treated in this
work.

A precedent worth recording comes from the second report's reading list: Giannoni,
Piccione and Verderesi study the stationary problem with
\(-\braket{\dot z, Y} \equiv k >0\) and restrict at once to \(|Y|<k\), which is
exactly Lemma~2. The restriction has stood since 1997, but what is studied here is
the new regime in which \(W\) is not Killing, when the boundary moves with
\(A(\eta)\).

\subsection{Major comment 2. Pontryagin extremals are not automatically local
minima.}
\textbf{Reply.} You are right. The error was more extensive than estimated: the
statement appeared in both versions of Paper~I. Both now report the hierarchy as
you wrote it, namely
\[
  \text{minimiser}\Longrightarrow\text{PMP extremal},
  \qquad
  \text{PMP extremal}\nRightarrow\text{local minimiser},
\]
with the note that local minimality requires a second variation or a
conjugate-point argument.

The revision reports sufficient conditions in the regime where they hold. On the
frozen arrival-time branch the optical geometry is Riemannian and the Gauss
curvature is negative throughout the exterior exactly when \(\hat{E}^2 \geq 3/2\).
There are then no equatorial conjugate points, and by the index theorem in the form
of Giannoni, Masiello and Piccione (\emph{J. Geom. Phys.} \textbf{35} (2000) 1--34)
every exterior extremal is a local minimiser. The surface is complete, so
Cartan--Hadamard on the universal cover returns one global minimiser for each
winding class, while the optical distance saturates the eikonal equation and
supplies the value function mentioned in your list.

These statements are equatorial by nature, not by present limitation, and the
manuscript says so. The tangential curvature is \(K_{\rm tan} = 1/B>0\) at the
circular orbit, so Cartan--Hadamard cannot hold in three spatial dimensions, and
the winding classes themselves do not survive the lift, since the equatorial
surface is \((2M,\infty)\times S^1\) while the full one is simply connected.

The out-of-plane Jacobi field is \(\sqrt{B}\,|\sin \varphi|\), which vanishes at a
swept azimuth of \(\pi\); for a fixed non-radial arc, \(K_t <0\) together with
\(\Phi < \pi\) is a sufficient condition for there to be no conjugate point in any
direction, and nothing stronger follows. In particular I do not import the bound
\(\Phi < \pi\) from Paper~I, which is proved there for the proper-time family,
whereas these statements live on the arrival-time family. Nothing analogous is
claimed on the proper-time family, where \(r^3K_\tau\to +M >0\) and the surface is
incomplete at the stationary limit.

The passage you were referring to, on perturbing the computed brachistochrones, is
no longer there; the appendix now states that such perturbations are a consistency
check and cannot certify a minimum in an infinite-dimensional class.

\subsection{Major comment 3. Endpoint protocol and costate relabelling.}
\textbf{Reply.} Accepted. Protocol~2 fixes the default launch, with a fixed spatial
target and a free arrival clock, the shooting being branch-dependent. Comparisons
between the two branches are made at \(J_t=AJ_{\eta}\), so that comparisons at the
same \(J\) and comparisons at the same endpoints are two different comparisons.

Section~3.4, the one place where the default does not apply, now names its two
protocols, records that at a common numerical value of \(J\) they compare different
orbits, and adds that an ordering between two selected extremals does not show that
either is the global minimiser. The algebraic regularity of the printed
Hamiltonians at \(f = 0\) is kept separate from the control-domain question.

\subsection{Major comment 4. The retrograde marginal orbit: cusp or corner.}
\textbf{Reply.} Accepted. This was the most useful comment in the report. I
rederived the case as you instructed, and both statements I had printed in the
paper were wrong.

Two exact identities organise it: \(p_r^2 = \mathcal{S}/(D^2_E\Delta ^2)\), with the
prefactor regular and nonzero at \(r_e\), and, at the marginal value, the exact
factorisation \(r\Delta -J_c^2D_E = (r-2M)(r^2+J_c^2)\), giving
\(\mathcal{S} = r^4f^2w(r^2+J_c^2)\). Hence \(r_e\) is a double root for both signs,
and it is not the reality condition that selects the prograde value. The
counter-rotation argument is withdrawn, and it was wrong in detail too, since near
\(r_e\) the retrograde marginal has \(\dot \varphi > 0\) and is locally co-rotating.

What breaks the \(\pm J\) symmetry is invisible in the radical but manifest in the
Hamilton equations: \(\tilde{p}_{\varphi}(r_e) = J-J_c\) is linear in \(J\), and
\(\mathcal{T}(r_e) = -a(J-J_c)/\bar{P}_e\) vanishes only at \(J=+J_c\). The
classification consists of four cases: a smooth periapsis for \(|J|>J_c\); a
semicubical cusp for \(|J|<J_c\); at \(J=-J_c\) a double root with \(r_e\) reached
only asymptotically, finite tangent and exponent~\(1\); and finally a regular
crossing at \(J=+J_c\).

The cusp in the main text was wrong, and the corner in the appendix was half right,
since the tangent is finite but there is no reflection. Your caution against merging
the generic and the marginal case is quantitatively justified: the cusp amplitude
\(K = (EJ/a^2)\sqrt{2M/(a^2-E^2J^2)}\) diverges as \(|J|\to J_c^-\).

The one-sided expansions, the signs, the branch orientation and the convergence
tests are in \texttt{verify\_cusp\_corner.wls} (21 checks) and in
\texttt{verify\_marginal\_hamilton.wls} (19 checks), all passing. For
\(M=1\), \(a=0.9\), \(E = 1.2\) the retrograde rate is \(C = -0.267000117\),
consistent with the closed form to \(10^{-6}\); the elapsed parameter increases by
\(34.4956\) per four decades of \(r-r_e\), against \(\ln(10^4)/|C| = 34.4955\).

All the material involved has been regenerated, not merely supplied with new
captions: the atlas had drawn a reflection at \(-J_c\), and Table~A1 is now indexed
on the multiplicity of the root and on whether it is attained.

\subsection{Major comment 5. The exterior separatrix as the principal result.}
\textbf{Reply.} Accepted, and adopted as the organising result. Proposition~10 gives
\(r_d = 3.513905124011657\,M\) and \(J_c^-=-8.053516003877019\), reproducing your
values to sixteen digits, with the printed closed form agreeing to \(10^{-12}\). At
\(r_d\) both control-domain inequalities hold strictly, \(g(W,W) = -0.430833\) and
\(\bar v^2 = 0.700811\), and they hold on the whole approach arc, so it is of
type~(E) rather than a continuation. A spin scan gives
\(r_d/M = 2.857675,\,3.204756,\,3.513905,\,3.579494\) for
\(a = 0.1,\,0.5,\,0.9,\,0.99\), so the exterior character is structural. The
prograde degeneration at \(r_d = 1.512292\,M\) lies outside \(r_+ = 1.435890\,M\)
but inside \(2M\), and is recorded as type~(A), not a physical control separatrix.
This also answers M4.

\subsection{Major comment 6. The tiered proof status of the fixed-endpoint result.}
\textbf{Reply.} Accepted. The conclusions distinguish five levels: an analytic
inequality for the integrand; a theorem for \(r_0\leq R_*(E,a)\); an asymptotic
result for \(r_0 \to \infty\) in the static case; a computer-assisted proposition on
a stated grid; and a conjecture in general. Protocol~4 assigns them different
evidential status. The text also says that an ordering between two selected
extremals does not show that either is the global minimiser. On executability you
were right that the CAP runner imported a module that was absent; see major
comment~10.

\subsection{Major comment 7. Scope and provenance of the first-order result.}
\textbf{Reply.} Accepted, and the problem ran deeper than the discrepancy you had
brought to light: the numbers had been transcribed by hand from different runs, and
re-running returned \(2.12\pm0.03\) and \(2.15\), so neither printed value was
current. They are now generated. One archived command,
\texttt{python3 paper2/provenance/make\_provenance.py}, runs each validation script
as a subprocess, reads the \(\varepsilon\)-window from the source, computes the hash
of every script, and writes the machine-readable record, the \LaTeX{} macros
included by the manuscript, and the manifest.

The theorem is now local on a compact regular subarc, excluding turning points, the
freezing surface, the stationary limit, coordinate singularities and crossings of
moving separatrices, the last identified as a singular inner problem.

As for the numbers: since the residual is \(c_2\varepsilon^2(1+k\varepsilon + \dots)\)
the local exponent is \(2+O(\varepsilon)\), so a fit over a finite window is biased
upward, and \(\sigma\) measures the quality of the fit rather than that bias.
Refitting over progressively smaller windows we obtain
\(2.118\to2.068\to2.043\), \(2.123\to2.072\to2.045\) and \(2.122\to2.070\to2.044\)
in the three configurations, each decreasing toward \(2\). The window is
\(\varepsilon \in \{0.001,0.002,0.004,0.008,0.016\}\) and the fitted interval is
\(r/M \in [8,11]\), at a distance of \(6M\) from the nearest excluded locus.

\subsection{Major comment 8. Higher-genus terminology.}
\textbf{Reply.} Accepted. A convention now establishes that ``length-two iterated
Abelian integral on the genus-two hyperelliptic curve'' is the only classification
claimed, and that ``genus-two dilogarithm'' is a descriptive shorthand for it, not a
claim that a function class with specified kernels, relations, monodromy or
single-valued completion has been defined. The three unqualified theorem-level uses
have been changed. ``Elliptic dilogarithm'' is retained, since at genus one the
object is established.

\subsection{Major comment 9. Physical language.}
\textbf{Reply.} Accepted. Protocol~3 fixes the dictionary and is applied
throughout: \(\hat E\) is the actively maintained rail charge, not a conserved
energy; \(J = p_\varphi\) is the axial control costate, not mechanical angular
momentum; the rail minimises a selected clock, not thrust or fuel; \(r_e=2M\) is the
conformal stationary limit; \(r_+\) is the null surface of the seed Kerr metric,
with the note that for a time-dependent \(A(\eta)\) neither horizon has been
computed; and inside \(2M\) we write ``continued orbit'', ``root pattern'' or
``analytic branch''.

\subsection{Major comment 10. Reproducibility package.}
\textbf{Reply.} All three findings are correct and I am glad they were checked.
\texttt{cap\_full.py} had been removed during development and never archived; it has
been restored in full from the prototype, with a self-test.

For the neighbouring directories the diagnosis turned out to be different: all
thirteen generators were tracked, but each inserted the repository root on
\texttt{sys.path} while the shared style module lives one directory further down. On
a fresh clone, therefore, all of them failed at import. Having corrected this
behaviour in all thirteen, a silent bug came to light: a helper function for the
turning radius returned the grid point just inside the forbidden region, dropping a
curve from a panel without raising any error.

The printed hashes were stale; the same two files carried three different values,
which is the signature of hand transcription. They are now generated by the command
given above.

Two statements of scope belong with this. The content digest covers the adiabatic
validation scripts, the windows and the fitted data listed in it; it is not a
checksum of the whole archive, and each manifest row separately carries the checksum
of its own file. The manifest associates the principal computational claims, and the
figures and tables that depend on them, not every equation.

\subsection{Major comment 11. Dependence on Paper I.}
\textbf{Reply.} Accepted. The incorrect sentence on the PMP is corrected there.
Paper~I's existence and normality theorem now defines the extended admissible class
and the clock state, states the regularity assumptions, explains the closure of the
non-autonomous clock equation, and identifies the endpoint protocol and the excluded
degeneracies.

The Vaidya qualification is carried into the abstract and the conclusions: the
compact control domain terminates at \(r=2m(v)\), and the calculations establish an
exterior approach and a contact threshold rather than an interior optimal
trajectory, while \(m'(v)<0\) in the ingoing metric is a formal continuation,
physical evaporation requiring the outgoing problem. Paper~II no longer imports
optimality claims from Paper~I; it cites the verification criterion as a criterion.

\subsection{Major comment 12. Relation to established work.}
\textbf{Reply.} Accepted and treated as a question of positioning.

A new paragraph in the introduction separates the established results from the new
ones. Established: Perlick's fixed-energy optical metric; the sub-Riemannian
formulation of Giannoni, Piccione and Verderesi; the arrival-time theory of Giannoni
and Piccione; the Morse theory of Giannoni, Piccione and Tausk. I have added the
Giannoni--Masiello--Piccione programme, including the finiteness results on
conformally stationary spacetimes, which is the published setting closest to the
frozen-\(A\) approach adopted in this work, and the work of Caponio, Corona, Giamb\`o
and Piccione on fixed-energy solutions of an indefinite Lagrangian with an affine
Noether charge.

The distinction is made explicit there: in all of those works the charge is
conserved, because the generator is a symmetry; here it is held fixed by control
against a background spacetime of which \(W\) is in general not a symmetry, at the
cost of proper acceleration, so that \(\hat E\) is a design parameter. Three
consequences have no stationary counterpart: the possible loss of compactness and
hence a causal boundary; the reduced Hamiltonian ceasing to be an interior first
integral; and the first-order response being a genuine perturbation problem.

I have also adopted the framing suggested in the closing part of the report, placing
the emphasis on the controlled invariant and on the causal boundary rather than on
one further exact trajectory.

\subsection{Minor comments.}
\begin{itemize}
    \item \textbf{M1} Done: Protocol~2 is stated once and used consistently.
    \item \textbf{M2} Done: \(J=p_\varphi\) is identified as a control costate at
    first use.
    \item \textbf{M3} Done: the compact boundary of the indicatrix and the convex
    body whose support function is used are distinguished, with an explanation of
    why no convexification is needed.
    \item \textbf{M4} Done: see major comment~5.
    \item \textbf{M5} Done. A mechanical check of all captions, re-run on the final
    text, found five that reported radii or momenta without the normalisation; all
    now include it.
    \item \textbf{M6} Done: every figure containing a curve inside \(r_e\) draws
    that portion in a different line style and states, in the legend or in the panel
    title, that the continuation is analytic and that no optimality is claimed.
    Figures whose curves terminate at \(r_e\) carry no such label, because there is
    no continuation in that case.
    \item \textbf{M7} Done: Protocol~3 carries a table keeping
    \(E\), \(\hat E\), \(E_{\rm eff}\), \(J\) and \(J_{\rm eff}\) distinct,
    including \(J_t = AJ_\eta\).
    \item \textbf{M8} Done: deviations are declared absolute unless otherwise
    indicated; norms are specified for vector residuals; every fitted slope is
    accompanied by its interval and, for the adiabatic convergence, by its
    \(\varepsilon\)-window.
    \item \textbf{M9} Done: Protocol~4 separates the levels of evidence and the
    claims are labelled.
    \item \textbf{M10} Done: the conclusion is shortened and lists only those
    results whose domain and proof status are established in the body of the paper.
\end{itemize}

\section{Referee 2}

\subsection{Bibliography.}
\textbf{Reply.} Accepted with thanks; the list was useful. I have added or made
explicit Giannoni, Piccione and Verderesi (1997); Giannoni and Piccione on
arrival-time brachistochrones; Giannoni, Piccione and Tausk, with the identifier
\texttt{math-ph/9905007} so that the preprint and the published version are visibly
the same work; Haws and Kiser (1995); and Ta\c{s} (\texttt{arXiv:2512.08776}).

Your suggestion of connecting the global analysis to the Morse index has shaped the
revision: it is the framework in which the extremal-versus-minimiser request is
properly answered.

One bibliographic note, so that the record is right: the article
\emph{J.\ Math.\ Phys.} \textbf{38}(12), 6367--6381 (1997), given as item~[6], was
written by Giannoni, Piccione and Verderesi. Giannoni, Masiello and Piccione are a
separate collaboration, and I now cite three of their papers in their own right,
including the finiteness results on conformally stationary spacetimes, which are the
published setting closest to the frozen-\(A\) regime used here. Checking your list
also brought to light a wrong page number in one of my own entries.

\subsection{Eigenvalues of the curvature operator and weighted manifolds.}
\textbf{Reply.} This comment sent me back to reconsider the geometry, and it was
worth it: there is a precise structural relation between the two frameworks which the
manuscript had left implicit, and a new section~2.2 now sets it out.

The relation is not the one the comment proposes, and the difference is one of
scope rather than of correctness; the theorem holds where its hypotheses are
satisfied, while the objects studied here fall outside them.

There are three connections in the paper now thanks to this report. First, the
optical reduction is a submersion, and with Perlick's metric on the base it is
horizontally conformal with dilation \(\Lambda^2 = \hat E^2 / [f(\hat E^2-f)]\), the
Perlick factor itself. Second, since \(\mathcal{L}_Wg=(2A'/A)g\), the selector is a
conformal Killing field and
\(\mathrm{d} \hat E/\mathrm{d}\tau|_{\text{geodesic}}=A'/A=\varepsilon\),
independently of the four-velocity: the rate at which an \emph{uncontrolled}
worldline drifts away from the rail is the adiabatic expansion rate, which makes
quantitative the middle rung of the
Killing\,\(\to\)\,conformal-Killing\,\(\to\)\,Kodama hierarchy. Third, choosing the
weight \(e^{-2f^2}=A^{-2}\) makes the weighted metric exactly the Kerr seed, with
\(\operatorname{div}_{\bar g}W = 0\) so that your equation~(32) holds verbatim,
while \(\operatorname{div}_{g}W = 4A'/A\); your Theorem~T8 requires a parallel
weight, and ours is parallel precisely at frozen \(A\). That framework therefore
describes exactly the frozen level of this construction, and the instinct behind the
comment is well founded.

In the revision I corrected my own statement of the first point. What I had called
the ``Fuglede--Ishihara dichotomy'' is Fuglede's definition of semiconformality
(\emph{Ann.\ Inst.\ Fourier} \textbf{28} (1978) 107, \S5), not a theorem, and
invoking it would have been circular. Following Ishihara's version instead gave a
definite answer.

His reduction of harmonicity to minimal fibres holds at unit dilation, whereas our
case has non-constant dilation, so the tension field carries a second term,
\(\tau(\pi) = -\,\mathrm{d}\pi(\mu+\nabla^H\ln\Lambda)\), and the non-geodesic
character of the fibres alone leaves open the possibility that the two terms cancel.
Computing both, with \(\mu = (M/r^2)\partial_r\), returns
\begin{align*}
    \mu + \nabla^H\ln \Lambda
      = \frac{M}{r^2}\,\frac{f}{\hat E^2-f}\,\partial_r\neq 0
      \qquad (r> 2M,\ \hat E > 1),
\end{align*}
strictly positive throughout the exterior. So the projection is not a harmonic
morphism, and the obstruction coincides with the quantity that makes the rail cost
something. The claim concerns the non-vanishing of a tension field computed on the
horizontal distribution, which is spacelike; I do not invoke the global elliptic
theory, which is Riemannian. The identity is verified in two independent
computer-algebra systems.

As for Myers, my first argument was too quick. Decay at infinity does not exclude
his hypothesis, and he himself supplies a counterexample. Tested on
\(\mathrm{Ric}(v,v)\) instead, the optical scalar curvature is negative at every
exterior point and for every rail energy, whereas \(\mathrm{Ric}(v,v)\ge e^2 > 0\)
would force \(R_{\text{opt}} \ge 3e^2 > 0\). The hypothesis fails by sign,
pointwise, and the relevant instrument is his Lemma, which requires no completeness.
Correcting a non sequitur produced a better result.

Set beside the objects treated here, Theorem~T2 concerns something else, for three
concrete reasons. Equation~(9) is a Hermite-type reduction of Abelian differentials
on a hyperelliptic curve; a dichotomy about the differential of a map and a
reduction of Abelian integrals are statements about different objects. A
controlled-rail brachistochrone is not a geodesic of a weighted metric, but an
extremal of a time-optimal control problem with an actively imposed pointwise
constraint, a compact indicatrix and a support Hamiltonian; there is moreover a
regime in which the control set loses compactness and the problem ceases to be well
posed, with no counterpart for geodesics. Finally, Remark~R4 stresses that the
photon follows geodesics, while Theorem~T25 constructs a Killing field from the
four-velocity; the rail is neither, and it is for this reason that the three-rung
hierarchy is needed.

For these reasons the results cited above do not follow from Theorem~T2. I should
like to make clear how narrow this conclusion is: it is not a judgement on the
theorem, nor on the framework, which I have adopted and cited. Three of the
connections this report proposed are now in the manuscript, together with the Myers
correction. The fourth, the implication from T2, is the one I have not managed to
make. Should it exist, I would be glad to cite it.

If the editor considers it useful to record the connection, I will add a neutral
remark to the effect that weighted-manifold and curvature-operator methods provide
an alternative setting for related variational problems, without asserting a
dependence.

\subsection{Experimental validation.}
\textbf{Reply.} Accepted, and the revision answers the request rather than declining
it. No gravitational validation is available and none is claimed, but the request
concerned techniques, and these exist for the parts of this work that do not involve
the Einstein dynamics. Appendix~A presents three levels.

The controlled rail is realisable on a bench: a holonomic planar robot or an \(XY\)
stage, tracked from above and commanded so that the \emph{measured} velocity
satisfies \(\dot x = c + R\,e(\theta)\). Applying 32--64 control directions on a grid
of states and fitting the velocity locus measures the indicatrix directly, comparing
it with the \(c\) and \(R\) predicted here, with no refitting of the trajectories; a
programmed \(A(\chi)\) then tests the first-order response and the
\(O(\varepsilon^2)\) scaling, and the bench can be driven across \(J =-J_c\) to
observe the four-case behaviour. This validates the reduced system and its control
structure, not the field equations, and I do not present it otherwise. Included in
the archive, under \texttt{Experimental/}, is a design study runnable with a single
command: the indicatrix identity closes to \(4.4 \times 10^{-16}\) and the
frozen-flow Hamiltonian is conserved to \(\sim 10^{-15}\). It is a design study, not
a set of data, and is labelled as such.

Analogue platforms supply the geometry. Rotational superradiance has been measured
in a draining vortex, with amplification \(14\%\pm 8\%\) at \(3.70\)~Hz,
\(R_{m=1} = 1.09\pm 0.03\) and \(R_{m=2} = 1.14\pm 0.08\) (Torres \emph{et al.},
\emph{Nature Phys.} \textbf{13} (2017) 833), and the same structure has since been
realised in superfluid \(^4\)He (\v{S}van\v{c}ara \emph{et al.}, \emph{Nature}
\textbf{628} (2024) 66). Their effective metric is a threading form with wind
\(\boldsymbol v\), hence of Randers type, and the ergosurface lies at
\(|\boldsymbol{v}| = c\), that is at \(g(W,W) = 0\), where the selector ceases to be
timelike: the first hypothesis of Lemma~2, and the same condition that terminates
our control domain at \(r = 2M\). It does not coincide with the energy-dependent
freezing condition \(\hat E = |W|\), at which \(W\) is still timelike.

The time-dependent component has its own realisations, in expanding condensates
(Eckel \emph{et al.}, \emph{Phys.\ Rev.\ X} \textbf{8} (2018) 021021) and in the
programmable scale factor of Viermann \emph{et al.} (\emph{Nature} \textbf{611}
(2022) 260). The two features of \(g = A(\eta)^2g_{\text{Kerr}}\) each have a
laboratory realisation and none that I know of contains both; combining them is the
experimental question the geometry poses.

Two limitations must be stated. These platforms carry waves, and therefore probe the
geometry and its boundary, not an actively controlled massive worldline; moreover, a
purely conformal rescaling is invisible to ideal null rays, so an analogue system
observes \(A(\eta)\) through the redshift of modes and particle production, not
through ray trajectories. It is because the rail is timelike at fixed \(\hat E\)
that the conformal factor acts on it at all, which is the same reason the problem
does not reduce to Fermat's principle.

Since no measurement is on offer, the revision states in one place the values an
independent recomputation must return: the separatrix at
\(r_d/M = 3.513905124011657\) with \(J_c^- = -8.053516003877019\), the threshold
\(\hat E^2 = 3/2\), the negativity of the optical scalar curvature, the twist
\(-a/2M\), the fibre acceleration, the minimum thrust and the adiabatic exponents.
Each is given with enough digits that a sign error cannot hide, and each localises a
different part of the construction. This is the form of falsifiability the problem
admits, and I prefer to commit to it rather than to assert an observational contact
the work does not have.

\section{A note of thanks}

Referee~1's report led to the identification of four errors: the cusp/corner
contradiction, the PMP statement, the missing module and the drifting hashes.
Pursuing them brought three more to light that the report could not have seen: an
incorrect import path that prevented every figure script from running on a fresh
clone; a silent error in the turning radius that removed a curve from a published
panel; and a wrong page number.

Referee~2's report is responsible for the material that is new rather than
corrected: the reading list, which made it possible to identify that page number and
the finiteness results; the insistence on experimental validation, which produced the
three-level section; and the geometric framework, which generated three connections
now present in the paper, together with the Myers correction. Where I have not been
able to follow that report, I have said so above and set out my reasons.

The manuscript is substantially better thanks to both.

\end{document}
